\begin{abstract}
We construct the field and area law of the unit-area critical Liouville
quantum sphere as a limit of ordinary subcritical quantum spheres, and
equip it with its critical intrinsic metric. We then prove that spherical
Fortuin--Kasteleyn planar maps at $q=4$, embedded by their equilateral flag
uniformizations, converge jointly to this sphere decorated by an independent
nested conformal loop ensemble $\CLE_4$. With deterministic distance
normalization, the convergence includes the area measure, the full embedded
distance function, and every macroscopic interface through all positive
integer edge counts.
\end{abstract}

\section{Introduction and statements}\label{sec:intro}

Random planar maps provide discrete models of random two-dimensional geometry. Decorating a map by a critical Fortuin--Kasteleyn random-cluster configuration \cite{FortuinKasteleyn1972} couples that geometry to a loop ensemble. The expected continuum objects are Liouville quantum gravity surfaces decorated by conformal loop ensembles. The endpoint $q=4$ is distinguished by critical quantum gravity and logarithmic corrections in the discrete encoding.

Sheffield's inventory model established a correspondence between decorated
planar maps and two-type inventory trajectories, and identified the transition
at $q=4$ \cite{SheffieldInventory}.  Duplantier, Miller, and Sheffield
constructed subcritical quantum surfaces by mating continuum trees
\cite{DMS}.  At the critical endpoint, Aru, Holden, Powell, and Sun established
the Brownian half-plane-excursion coupling for quantum disks and the
near-critical convergence of their fields and quantum measures
\cite{AHPScritical}.  The limiting Brownian data do not by themselves determine
the decorated critical surface.  Da Silva, Hu, Powell, and Wong proved the
critical infinite-volume FK encoding limit with its logarithmic discrepancy normalization and
an exact fully packed loop description \cite{DSHPW}.  Our geometric theorem
retains the actual conformal structure, the original graph metric, area, and
all macroscopic interfaces in one finite-map limit.

The area construction builds on critical Gaussian multiplicative chaos.
Duplantier--Rhodes--Sheffield--Vargas constructed the derivative chaos and
identified its relation to Seneta--Heyde normalization
\cite{DRSVderivative,DRSV}.  Powell proved universality of derivative-normalized
convolution approximations under stated mollifier hypotheses \cite{Powell};
the circle-average Seneta--Heyde construction is treated by
Huang--Rhodes--Vargas \cite[Theorem~4.2]{HRVcriticalDisk}.
Aru--Powell--Sep\'ulveda proved the fixed-field limit
$(2-\gamma)^{-1}\mu^\gamma\to2\mu'$ \cite[Theorem~1.1]{APS}.
Section~\ref{sec:continuum} specializes the corrected cutoff argument to
circle averages and passes to the varying, area-normalized sphere law.
Its starting law is the ordinary subcritical sphere of
\cite[Definition~4.21(ii) and Proposition~A.13]{DMS}.

For $0<\gamma<2$, Aru--Huang--Sun identify the DMS unit-area sphere
with three independent area marks with the unit-volume sphere of
David--Kupiainen--Rhodes--Vargas with three insertions of weight $\gamma$
\cite[Theorem~1.1]{AHSSphereEquivalence}.
The latter authors also discuss sphere constructions at
$\gamma=2$ \cite[Section~4.1]{DKRVSphere}.
Here we study the critical limit of the ordinary DMS area-one law,
retaining the area-sampling interpretation of its marks.

The metric construction uses the critical local metric and Euclidean topology
of Ding--Gwynne \cite{DGuniqueness,DGtopology}.  We prove the additional
critical coordinate-change and sphere-completion statements in
Section~\ref{sec:continuum}; the coordinate argument follows the
local linearization and annular comparison strategy of Gwynne--Miller
\cite[Sections~2--3]{GMCovariance}, using the Euclidean-topology estimates
available at criticality.  The later rigidity argument adapts Ding--Gwynne's
optimal-comparison and field-perturbation strategy to discrete passage
profiles, for which the required conditional laws must first be established.
Field reconstruction uses the log-mass strategy of
Berestycki--Sheffield--Sun and the critical reconstruction framework of Vihko
\cite{BSSReconstruction,VihkoReconstruction}, with the local normalization and
comparison-class argument given in Section~\ref{sec:conformal}.

Quantum-disk exploration provides a separate source of geometric information.
Miller--Sheffield--Werner established the subcritical simple-CLE disk cutting
laws \cite{MSWQuantum}, while Kammerer identified the critical total-perimeter
process and related exploration distances \cite{Kammerer}.
Ang--Gwynne proved the conditional independent-disk law for components of an
independent $\CLE_4$ on a critical disk and identified their boundary lengths
through a reweighted $3/2$-stable process
\cite[Theorems~1.1--1.2]{AngGwynneCritical}.
We retain these antecedents while proving the marked two-side cutting law,
exact attachment rules, and finite-map comparisons needed here.
The conformal loop and field inputs come from the loop-soup construction and
sphere nesting \cite{SWcle,KWnested}, the GFF contour coupling
\cite{SchrammSheffieldContour}, and the subsequent signed exploration,
level-line, and bounded-type local-set theories
\cite{MSWTrunk,WangWu,ASWlocal,ASlocal}.

The two main conclusions are the independently specified critical sphere of
Theorem~\ref{thm:continuum} and the joint finite-map limit of
Theorem~\ref{thm:main}. The passage from the FK encoding to this geometric
limit has three distinct stages. Exact contacts and incidence comparisons
first identify the entire flag sphere up to a homeomorphism. Conditional
laws for local observations then identify its prescribed conformal
uniformization and compare its graph distances with the critical metric.
Finally, an estimate controlling every vertex extends the local distance
comparison to the full embedded distance function. Each stage is proved
at $q=4$; no convergence theorem for FK maps at another value of $q$ is used.

\subsection{The finite model and its canonical embedding}

Let $\mathcal F_n$ be the set of pairs $(M,A)$ in which $M$ is a rooted connected planar multigraph with $n$ edges in the oriented sphere and $A\subseteq E(M)$. Maps are considered up to orientation-preserving homeomorphism preserving the distinguished oriented root edge and $A$. Loops and multiple edges are allowed. The spanning subgraph $(V(M),A)$ includes its isolated vertices. If $M^*$ is the dual map, let $A^*$ consist of the dual edges crossing $E(M)\setminus A$. Write $k$ for the number of components and set
\begin{equation}\label{intro:fk-weight}
 \ell(M,A)=k(A)+k(A^*)-1
 =2k(A)+|A|-|V(M)|,
 \qquad
 \Prob\{(M_n,A_n)=(M,A)\}
 =\frac{2^{\ell(M,A)}}{Z_n(4)}.
\end{equation}
Here $Z_n(4)$ is the sum of the displayed weights over $\mathcal F_n$. In particular, the factor depending on $|V(M)|$ is part of the law. Let $d_n$ be shortest-path distance on the original primal vertices, using every original edge with length one.

For each endpoint-and-side incidence occurrence of an original edge, take an equilateral Euclidean triangle with corners labelled $v,e,f$: the original vertex, the edge midpoint, and the face center. Repeated occurrences at loops or bridges remain distinct. Glue the $4n$ triangles according to their incidence identifications. The resulting oriented polyhedral sphere $X_n$ has its conformal structure extended across its conical vertices. Let $m_n$ be normalized Euclidean area, giving mass $1/(4n)$ to each triangle.

Conditionally on $X_n$, sample $x_n^1,x_n^2,x_n^3$ independently from $m_n$. Almost surely they are distinct. Let
\[
 \phi_n:X_n\longrightarrow\widehatC,
 \qquad \phi_n(x_n^1)=0,\quad
 \phi_n(x_n^2)=1,\quad \phi_n(x_n^3)=\infty,
\]
be the unique orientation-preserving conformal uniformization
\cite[Theorem~II.2.15]{SaintGervaisUniformization}. This prescription is part of the discrete model. Let $\rho$ denote the standard spherical distance on $\widehatC$.

We also fix the interface drawing. In a flag triangle retain the closed region whose barycentric coordinate at $v$ is at least $2/3$. If its original edge belongs to $A_n$, retain in addition the region whose barycentric coordinate at $f$ is at most $1/3$. Glue these regions with the flags. Their polygonal boundary components are the FK interfaces. Let $\Gamma_n$ be their images under $\phi_n$, as unrooted, unoriented loops, retaining every nesting level and every occurrence. Figure~\ref{fig:flags} depicts the local prescription. The flag geometry defines the conformal structure and drawing; $d_n$ is the original primal graph distance.

\begin{figure}[htbp]
\centering
\begin{tikzpicture}[scale=1.05,line join=round,font=\small]
\definecolor{flagblue}{RGB}{28,86,141}
\definecolor{flagfill}{RGB}{218,234,247}
\begin{scope}
\coordinate (v) at (0,0);\coordinate (e) at (3,0);\coordinate (f) at (1.5,2.598);
\fill[flagfill] (v)--(1,0)--(.5,.866)--cycle;
\draw[thick] (v)--(e)--(f)--cycle;
\draw[flagblue,very thick] (1,0)--(.5,.866);
\node[below left] at (v) {$v$};\node[below right] at (e) {$e$};\node[above] at (f) {$f$};
\node[below=14pt] at (1.5,0) {unoccupied edge};
\node at (1.5,1.35) {$\lambda_v\geq2/3$};
\draw[-{Stealth[length=1.5mm]},thin] (1.03,1.08)--(.42,.33);
\end{scope}
\begin{scope}[xshift=5.2cm]
\coordinate (v2) at (0,0);\coordinate (e2) at (3,0);\coordinate (f2) at (1.5,2.598);
\fill[flagfill] (0,0)--(3,0)--(2.5,.866)--(.5,.866)--cycle;
\draw[thick] (v2)--(e2)--(f2)--cycle;
\draw[flagblue,very thick] (.5,.866)--(2.5,.866);
\node[below left] at (v2) {$v$};\node[below right] at (e2) {$e$};\node[above] at (f2) {$f$};
\node[below=14pt] at (1.5,0) {occupied edge};
\node at (1.5,.36) {$\lambda_f\leq1/3$};
\end{scope}
\end{tikzpicture}

\caption{The retained region in one equilateral flag. The left panel applies when the original edge is unoccupied; the right panel adds the strip next to $ve$ when it is occupied. The blue segments inside a flag are interface pieces; sides shared with adjacent flags are glued before the boundary is taken. Labels specify incidence occurrences, so the same construction applies to loops and bridges.}
\label{fig:flags}
\end{figure}

The embedded distance graph at scale $a>0$ is the compact set
\begin{equation}\label{intro:distance-graph}
 K_n(a)=\{(\phi_n(u),\phi_n(v),a d_n(u,v)):
                  u,v\in V(M_n)\}.
\end{equation}

\subsection{The independently specified continuum law}

For $\sqrt2<\gamma<2$, take the ordinary unit-area $\gamma$-quantum sphere of \cite[Definition~4.21(ii)]{DMS}, disintegrated at area one, with its original marked points forgotten. Sample three independent points from its area measure, and use the orientation-preserving coordinate sending them to $0,1,\infty$. Write $h_\gamma$ for the field in that coordinate. We use the circle-average convention
\[
 \mu_h^\gamma(\dd z)
  =\lim_{\eps\downarrow0}
       \eps^{\gamma^2/2}e^{\gamma h_\eps(z)}\,\dd^2z.
\]
Define
\begin{align}
 Q_\gamma&=\frac2\gamma+\frac\gamma2,
 &c_\gamma&=\frac1\gamma\log\bigl(2(2-\gamma)\bigr),
 &\widehat h_\gamma&=h_\gamma+c_\gamma,\label{intro:renormalization}\\
 \nu_\gamma&=\frac{\mu_{\widehat h_\gamma}^\gamma}{2(2-\gamma)}
             =\mu_{h_\gamma}^\gamma,
 &H_\gamma&=\widehat h_\gamma-
                \frac{Q_\gamma}{2}\log g_{\rm round},
 &g_{\rm round}(z)&=\frac4{(1+|z|^2)^2}.
 \label{intro:round-field}
\end{align}
The expression in the chart at infinity agrees, so $H_\gamma$ is a scalar distribution on the round sphere. The factor in \eqref{intro:renormalization} is fixed by the local normalization
$\mu^\gamma/[2(2-\gamma)]\to\mu^{\rm crit}$ for a fixed ordinary GFF \cite{APS}; the convergence of the varying marked sphere laws is part of the following theorem.

\begin{theorem}[The critical area-marked sphere]\label{thm:continuum}
As $\gamma\uparrow2$, the laws of $(H_\gamma,\nu_\gamma)$ converge to a unique nondegenerate probability law $\mathsf P_{\rm crit}$ in
\[
 H^{-1}(S^2,g_{\rm round})\times\mathcal P(S^2),
\]
with the norm topology in the first factor and the weak topology in the second. Let $(H,\nu)$ have this law and set $h=H+\log g_{\rm round}$ in local conformal coordinates. Then $\nu=\mu_h^{\rm crit}$, where on compact sets avoiding $0,1,\infty$,
\begin{equation}\label{intro:critical-area}
 \mu_h^{\rm crit}(\dd z)
 =\lim_{\eps\downarrow0}
  \bigl(2\log(1/\eps)-h_\eps(z)\bigr)
        \eps^2e^{2h_\eps(z)}\,\dd^2z
\end{equation}
locally weakly in probability. This positive measure extends across the three marks without atoms, has full support, and has total mass one.

After forgetting the three marks, sampling three new conditionally independent area points and normalizing them to $0,1,\infty$ recovers $\mathsf P_{\rm crit}$. If $f$ maps the new coordinate to the old one, the field transforms as
\[
 h_{\rm new}=h_{\rm old}\circ f+2\log|f'|.
\]

With any one fixed deterministic normalization of the Ding--Gwynne critical local metric, the compatible chart metrics define an intrinsic length metric on the sphere punctured at the three marks. Its completion adds exactly those marks and is a finite continuous metric $D_h$ on $S^2$, inducing the usual topology. The area and completed metric are measurable functions of the same field and obey the corresponding charge-$2$ coordinate rules.
\end{theorem}

Conditionally on $(H,\nu)$, let $\Gamma$ have the M\"obius-invariant whole-plane nested $\CLE_4$ law on the sphere \cite{KWnested}, independently of the entire marked surface. This fixes the decorated continuum target. No convergence of subcritical metrics as $\gamma\uparrow2$ is part of Theorem~\ref{thm:continuum} or needed in its proof.

\subsection{Joint convergence}

We first specify the loop topology. For continuous loops $\eta,\eta':S^1\to S^2$, let $d_{\rm loop}(\eta,\eta')$ be the infimum of their uniform spherical distance over circular orientation-preserving reparametrizations and reversal. Identify representatives at zero distance. A loop collection is locally finite if it contains only finitely many loops of diameter above any positive threshold. An $\eps$-matching of two such collections is a partial bijection that covers every loop of diameter greater than $\eps$ on each side, matches loops at $d_{\rm loop}$-distance at most $\eps$, and leaves only loops of diameter at most $\eps$ unmatched. Loop convergence means the existence of such matchings with $\eps\to0$.

For compact subsets of $S^2\times S^2\times[0,\infty)$, use Hausdorff distance for the product metric
\[
 \rho(z,z')+\rho(w,w')+|r-r'|.
\]

\begin{theorem}[Joint critical universality]\label{thm:main}
For the finite law \eqref{intro:fk-weight}, there is a deterministic sequence $a_n>0$, with $a_n\to0$, such that through all positive integers $n$,
\begin{equation}\label{intro:joint-limit}
 \bigl((\phi_n)_*m_n,K_n(a_n),\Gamma_n\bigr)
 \ \xrightarrow{\ d\ }\
 \bigl(\mu_h^{\rm crit},
       \{(z,w,D_h(z,w)):z,w\in S^2\},\Gamma\bigr).
\end{equation}
Here $(H,\mu_h^{\rm crit})\sim\mathsf P_{\rm crit}$, the area topology is weak convergence, the distance-graph topology is the Hausdorff topology just defined, and the loop topology is convergence by $\eps$-matchings. All nesting and multiplicities are retained. In the same joint convergence,
\begin{equation}\label{intro:flag-mesh}
 \max_{T\ \text{flag of }X_n}\diam_\rho\phi_n(T)
       \ \xrightarrow{\ \Prob\ }0.
\end{equation}
The metric normalization is the one fixed in Theorem~\ref{thm:continuum}; changing it by a deterministic factor changes $a_n$ by that factor.
\end{theorem}

The Hausdorff assertion in \eqref{intro:joint-limit} controls every approaching pair of vertices. It therefore includes possible exceptional endpoints and does not require selecting a finite set of typical vertices. The theorem makes no claim of an explicit power law for $a_n$, a convergence rate, or a formula for the critical metric dimension.

Forgetting the embedding gives the following metric-measure consequence.
For a loop edge, each of its two endpoint occurrences contributes one to
the degree.

\begin{corollary}[The compact metric-measure limit]\label{intro:ghp}
Define a probability measure on the primal vertices by
\[
 \eta_n=\sum_{v\in V(M_n)}\frac{\deg_{M_n}(v)}{2n}\,\delta_v.
\]
With the deterministic normalization of Theorem~\ref{thm:main},
\[
 (V(M_n),a_nd_n,\eta_n)
 \ \xrightarrow{\ d\ }\
 (S^2,D_h,\mu_h^{\rm crit})
\]
in the Gromov--Hausdorff--Prokhorov topology, jointly with the embedded
area and loop convergence in that theorem.
\end{corollary}

\begin{proof}
This is a deterministic consequence of the distance-graph and flag-mesh
assertions of Theorem~\ref{thm:main}. Work on a representation on which
those assertions and the weak area convergence hold almost surely.
Every endpoint occurrence has two side flags, so sending each flag to
its primal vertex pushes $m_n$ to $\eta_n$, including at loops and
bridges. The spherical displacement of this projection is at most the
maximal flag mesh. Consequently
$\widetilde\eta_n=(\phi_n)_*\eta_n$ converges weakly to
$\mu_h^{\rm crit}$, and the embedded vertices become dense in $S^2$.

Continuity of $D_h$ and Hausdorff convergence of the full distance graphs
give
\[
 \sup_{u,v\in V(M_n)}
 \bigl|a_nd_n(u,v)-D_h(\phi_n(u),\phi_n(v))\bigr|
 \longrightarrow0.
\]
Indeed, a sequence violating this assertion has a convergent graph-point
subsequence, whose limit belongs to the graph of $D_h$.
Choose $r_n\downarrow0$ large enough to cover the maximal flag mesh and
to couple $\widetilde\eta_n$ with $\mu_h^{\rm crit}$ so that all but
$r_n$ of the mass has spherical displacement at most $r_n$. The relation
$\{(v,z):\rho(\phi_n(v),z)\le r_n\}$ covers both spaces. Its distortion
for $a_nd_n$ and $D_h$ tends to zero by the displayed bound and uniform
continuity of $D_h$, and it carries all but $r_n$ of the lifted coupling.
The correspondence-and-coupling characterization of
Gromov--Hausdorff--Prokhorov convergence proves the assertion. The same
construction on each coupled subsequence proves the stated joint
convergence in law.
\end{proof}

\begin{table}[htbp]
\centering\small
\begin{tabular}{@{}p{0.30\textwidth}p{0.65\textwidth}@{}}
\toprule
Quantity & Convention \\
\midrule
Model parameter & $q=4$ throughout the map theorem; critical LQG parameter $\gamma=2$. \\[3pt]
Discrete size and area & $n$ original edges, $2n$ Tutte triangles, $4n$ flags; each flag has mass $1/(4n)$. \\[3pt]
Graph metric & Every original primal edge has length one; the deterministic unit is $a_n$. \\[3pt]
Perimeter-count unit & $B_n=(2/\pi)\sqrt n\log n$, distinct from the graph-distance unit. \\[3pt]
Critical bulk area & Derivative normalization in \eqref{intro:critical-area}; near-critical divisor $2(2-\gamma)$. \\[3pt]
Conformal normalization & Three independent area samples, sent in order to $0,1,\infty$. \\
\bottomrule
\end{tabular}
\caption{Normalizations retained in both main theorems.  Boundary-length units for auxiliary disks are fixed separately in Section~\ref{sec:proxy}.}
\label{tab:normalizations}
\end{table}

\subsection{Proof structure and reusable arguments}

\paragraph{The target surface and its metric.}
Section~\ref{sec:continuum} proves Theorem~\ref{thm:continuum}
independently of the finite maps. A maximum-centered cylinder
representation exposes the area-conditioning tilt. Uniform tail estimates
permit passage to criticality, and local absolute continuity separates the
observed field from its area marks when the coordinates change. The metric
is then constructed directly from critical local metrics: annular
comparisons prove conformal covariance, and resampling the three marks
completes the metric at the original marks. The field limit and the metric
construction are separate arguments.

\paragraph{The topological comparison.}
Sections~\ref{sec:discrete}--\ref{sec:cutting} retain the exact finite
weight and its perimeter and volume units. A strictly subcritical boundary
proxy controls the images of small quantum arcs and children in the
critical limit. The resulting contact rule specifies exactly which
boundary occurrences represent the same point. Together with discrete
incidence estimates, it yields homeomorphisms comparing the entire flag
sphere, complete loop tours, and area. A loop-counting Palm identity and
unexamined volume reservoirs identify the ordinary area-one sphere law
after conditioning on the exact finite size. At this stage the
homeomorphisms provide topological coordinates; Section~\ref{sec:conformal}
will identify the prescribed flag uniformizations.

\paragraph{Conditional laws for local observations.}
Sections~\ref{sec:walls} and~\ref{sec:kernels} supply the common input
to the conformal and metric arguments. Auxiliary GFF level lines control
the possible local exploration orders when the exterior changes. Integer
gluing constraints, including their residues, describe which raw map
blocks can be assembled. An independent reference experiment is changed
by a density depending on the chart's macroscopic parameters, increment
paths, and exterior discrepancy variables. Once these data are fixed,
the density does not depend on the retained local payloads, which can
therefore be disintegrated over their actual continuation trees. This gives product
conditional laws for buffered observations even after the entire primary
tree is known. One finite choice then describes the whole countable vector
of observation kernels, including all later cost units. These exact
conditional laws make it possible to test many local candidates and to
compare them after a change of the field outside the observed patch.

\paragraph{The actual conformal coordinates.}
Section~\ref{sec:conformal} applies those laws to extremal length on the
flag surface. A transverse-flow argument prevents conformal collapse and
gives a quasiconformal limiting chart. Local determination makes its
distortion tensor deterministic, and the behavior under coordinate
rotations makes that tensor scalar. The three-point normalization therefore
identifies the limiting conformal map. Field reconstruction and elimination
of auxiliary exploration sources put the subsequent comparisons in the
law of the intended pair of fields.

\paragraph{Metric comparison on ordinary patches.}
Section~\ref{sec:metric} compares primal path costs directly with the
critical reference metric. If the optimal lower and upper comparison
factors differed, a fixed local field perturbation would create enough
strict shortcuts to contradict a nearly optimal path. The factors
therefore agree. This also supplies deterministic scale bounds. The
comparison at this point uses paths with open endpoint regions; it does
not yet attach every discrete vertex to its limiting point.

\paragraph{Every endpoint and one deterministic unit.}
Section~\ref{sec:endpoints} reduces that remaining attachment problem
to a bound on bad grid times in balanced peeling segments. Complete
incidence records ensure that a separating path found in an ordinary
local experiment remains valid in every compatible filling of the
unobserved map. Section~\ref{sec:address} counts the possible bad addresses
and combines equality constraints with survival estimates to prove the
required bound. The resulting uniform endpoint modulus identifies the
entire compact distance graph. Finally,
a bounded monotone calibration removes the deterministic scalar on every
subsequence, without distance-moment assumptions, and proves the full
sequence assertion of Theorem~\ref{thm:main}.

Figure~\ref{fig:proof-dependencies} records the logical order.  The
continuum field/area limit and critical metric construction are established
before the finite-map comparisons.  The strictly subcritical proxy in
Section~\ref{sec:proxy} only controls small quantum arcs and offspring; the
critical reference metric enters the separate distance-identification step.

\begin{figure}[p]
\centering
\begin{tikzpicture}[
 box/.style={draw,rounded corners=2pt,align=center,text width=5.6cm,
              minimum height=1.05cm,inner sep=6pt,font=\small},
 wide/.style={box,text width=8.9cm},
 arr/.style={-{Stealth[length=2mm]},semithick}]
\node[box] (sphere) at (-3.25,0)
 {Critical field, area, and marks\\\S\ref{sec:continuum}; Theorem~\ref{thm:continuum}};
\node[box] (finite) at (3.25,0)
 {Exact FK$(4)$ weights and disk volumes\\\S\ref{sec:discrete}};
\node[wide] (cut) at (0,-2.15)
 {Boundary proxy and marked critical cutting\\Exact boundary contacts and flag incidences\\
  \S\ref{sec:proxy}--\S\ref{sec:cutting}};
\node[wide] (walls) at (0,-4)
 {GFF walls and finite local exploration orders\\\S\ref{sec:walls}};
\node[wide] (kernels) at (0,-5.85)
 {Integer gluing, residues, and common conditional kernels\\\S\ref{sec:kernels}};
\node[wide] (conformal) at (0,-7.7)
 {Extremal length identifies the actual flag uniformizations\\\S\ref{sec:conformal}};
\node[wide] (metric) at (0,-9.55)
 {Open-port metric rigidity and deterministic scale bounds\\\S\ref{sec:metric}};
\node[wide] (end) at (0,-11.4)
 {Balanced peeling and the joint address bound\\
  Uniform control of every endpoint; \S\ref{sec:endpoints}--\S\ref{sec:address}};
\node[wide] (joint) at (0,-13.3)
 {Bounded calibration and the full joint limit\\
  Lemma~\ref{addr:calibration}; Theorem~\ref{thm:main}};
\draw[arr] (sphere.south) -- (cut.north west);
\draw[arr] (finite.south) -- (cut.north east);
\draw[arr] (cut) -- (walls);
\draw[arr] (walls) -- (kernels);
\draw[arr] (kernels) -- (conformal);
\draw[arr] (conformal) -- (metric);
\draw[arr] (metric) -- (end);
\draw[arr] (end) -- (joint);
\draw[arr] (sphere.west) -- ++(-.4,0) |- (metric.west);
\node[font=\scriptsize,rotate=90] at (-6.2,-5.2)
 {Critical metric charts and completion, \S\ref{sec:continuum}};
\end{tikzpicture}

\caption{Proof dependencies within this article.  The critical sphere
construction is independent of finite maps.  Exact cutting and finite local
kernels feed both conformal identification and metric rigidity.  The final
endpoint bound combines balanced peeling with the joint address estimate,
then bounded calibration fixes one deterministic unit through all map sizes.}
\label{fig:proof-dependencies}
\end{figure}

\subsection{Conventions}

All random fields and conditional distributions are realized on standard Borel spaces; conditional kernels are chosen measurably, with their versions fixed when introduced. On a non-Jordan boundary, an endpoint or boundary order refers to prime ends or to explicitly specified incidence occurrences. A protected buffer means a strict open spatial or label margin, whose size is fixed before the corresponding limit. Constants may depend on such fixed localizations. Perimeter count units, graph-cost units, and Euclidean coordinate scales are kept distinct. The notation used for auxiliary scales and cut regions is local to its section.
